\par\noindent\textbf{(1)} 横向位移xi≈lθ\_i，动能T=ml²(θ̇1²+θ̇2²)/2，重力势U\_g=mgl(θ1²+θ2²)/2，弹簧势U\_s=kl²(θ1−θ2)²/2。故矩阵方程θ̈+[(g/l+k/m),−k/m;−k/m,(g/l+k/m)]θ=0，质量矩阵为单位阵。
\par\noindent\textbf{(2)} 正交本征向量e+=(1,1)/√2、e−=(1,−1)/√2，对应ω+²=g/l、ω−²=g/l+2k/m。弹簧对同相模不伸缩，而反相相对位移增大，频率提高；两模对称正交，可独立传播任意初值。
\par\noindent\textbf{(3)} 初态θ1(0)=A、θ2(0)=0且初角速度都零，模态坐标初始均A/√2，故θ1=A[cos(ω+t)+cos(ω−t)]/2，θ2=A[cos(ω+t)−cos(ω−t)]/2。利用三角和差可表示载波(ω++ω−)/2与拍包络(ω−−ω+)/2。
\par\noindent\textbf{(4)} k/m≪g/l时频差Δω≈(k/m)/√(g/l)，初始能量的慢交换按sin²(Δω t/2)变化；第一次近乎全部到第二摆的时间t\_ex≈π/Δω，强耦合时两频率差不小，不宜将总瞬时能量直接等同单一慢包络平方。
